% !TEX program = lualatex
% Adapted from unimath-symbols and unimath-euler.ltx; see LICENSE.
\documentclass[a4paper,final]{article}
\makeatletter

\usepackage{ragged2e,setspace,booktabs,catchfile,shortvrb,geometry,graphicx,metalogo,textcomp,longtable,array,xcolor}
\geometry{margin=24mm}
\setlength{\parindent}{0pt}

\def\cmd#1{\texttt{\textbackslash\expandafter\@gobble\string#1}}

\usepackage{fontspec}
\setmainfont{texgyrepagella}%
 [
  Extension = .otf,
  UprightFont = *-regular,
  ItalicFont = *-italic,
  BoldFont = *-bold,
  BoldItalicFont = *-bolditalic,
  Ligatures=TeX,
  Numbers={Lowercase,Monospaced},
 ]
\usepackage[math-style=literal,bold-style=literal]{times-math}
\usepackage[hidelinks,pdftitle={Times Math glyph reference},pdfauthor={Junhao Wu}]{hyperref}

\ExplSyntaxOn
\cs_generate_variant:Nn \fontspec_set_fontface:NNnn {c}
\cs_generate_variant:Nn \fontspec_set_family:Nnn {cV}

\def\FINISH{}
\newcount\Nfonts
\def\defmathfont#1#2#3{
  \Nfonts=\numexpr\Nfonts+1\relax
  \tl_set:cn {MTcolour#1} {#3}
  \newcounter{#1}
  \setcounter{#1}{-1}
  \fontspec_set_fontface:cNnn {#1} \x
    {ItalicFont={},BoldFont={}} {#2}
  \tl_set:Nn \l_tmpa_tl {ItalicFont={},BoldFont={}}
  \sys_if_engine_luatex:T
    {\tl_put_right:Nn \l_tmpa_tl {,Renderer=HarfBuzz}}
  \fontspec_set_family:cVn {MTfamily#1} \l_tmpa_tl {#2}
  \tl_gput_right:Nn \FINISH {
    \refstepcounter{#1}\label{count:#1}
  }
}
\def\defrawmathfont#1#2#3{
  \defmathfont{#1}{#2}{#3}
  \fontspec_set_family:Nnn \l_tmpa_tl {Script=Math,Colour=#3} {#2}
  \exp_args:NnnV \DeclareSymbolFont {MT#1} {TU} \l_tmpa_tl {m} {n}
}
\def\MTRawOperator#1#2#3{
  \Umathchar 1~ \use:c {symMT#1}~ #2\scan_stop:
  \tl_if_in:NnT \l__um_nolimits_tl {#3} {\nolimits}
  \sb{0}\sp{1}
}
\def\MTRawOverlay#1#2#3{
  \sys_if_engine_luatex:TF
    {\Umathaccent overlay~fixed~0~ \use:c {symMT#1}~ #2\scan_stop: {\MTcircle}}
    {\Umathaccent 0~ \use:c {symMT#1}~ #2\scan_stop: {\MTcircle}}
}
\ExplSyntaxOff

\def\defmathfallback#1{
  \Nfonts=\numexpr\Nfonts+1\relax
  \newcounter{#1}
  \setcounter{#1}{-1}
  \g@addto@macro\FINISH{%
    \refstepcounter{#1}\label{count:#1}
  }
}

\defrawmathfont{xits}{XITSMath-Regular.otf}{CC6666}
\defrawmathfont{tgm}{texgyretermes-math.otf}{66CC66}
\IfFontExistsTF{Times New Roman}
  {\defmathfont{tnr}{Times New Roman}{6666CC}}
  {\defmathfallback{tnr}}
\defmathfont{tgt}{texgyretermes-regular.otf}{CC66CC}


\Nfonts=\numexpr\Nfonts+1\relax
\begingroup\ExplSyntaxOn
\newcounter{mtn}\setcounter{mtn}{-1}
\tl_gput_right:Nn \FINISH { \refstepcounter{mtn}\label{count:mtn} }
\ExplSyntaxOff\endgroup

\newif\ifMTok
\def\GateOne#1#2#3{%
  \expandafter\iffontchar\csname#1\endcsname #3\relax
    \MToktrue \else \MTokfalse \fi}
\def\MTbody#1{#1}
\def\MTcircle{\hbox{\scalebox{.65}{\hbox{\xits\char"25CC}}}}
\def\MTwidebase{\MTcircle\MTcircle\MTcircle}
\def\MTRawBody#1#2#3{\Umathchar 0 \csname symMT#1\endcsname #2\relax}
\def\MTRawMark#1#2{%
  \hbox{\fontfamily{\csname MTfamily#1\endcsname}\selectfont
    \ooalign{\hfil\MTcircle\hfil\crcr
      \hfil\char"00A0\char#2\relax\hfil\crcr}}}
\ExplSyntaxOn
\cs_new_protected:Npn \MTRawText #1#2
  {
    \bool_lazy_or:nnTF
      {\int_compare_p:n {"0300 <= #2 <= "036F}}
      {\int_compare_p:n {"20D0 <= #2 <= "20FF}}
      {\MTRawMark{#1}{#2}}
      {\hbox{\use:c {#1}\char#2\scan_stop:}}
  }
\ExplSyntaxOff
\def\SYMC#1#2#3#4{%
  #4{#1}{#3}%
  \ifMTok
    \refstepcounter{#1}%
    \begingroup
    \color[HTML]{\csname MTcolour#1\endcsname}%
    \ifcsname symMT#1\endcsname
      \Umathcode`x=0 \csname symMT#1\endcsname `x\relax
      \Umathcode`y=0 \csname symMT#1\endcsname `y\relax
      \Umathcode`+=2 \csname symMT#1\endcsname `+\relax
      $\displaystyle\MTRawBody{#1}{#3}{#2}$%
    \else
      \MTRawText{#1}{#3}%
    \fi
    \endgroup
  \fi}

\def\ColXITS#1#2{\SYMC{xits}{#1}{#2}{\GateOne{xits}}}
\def\ColTGM#1#2{\SYMC{tgm}{#1}{#2}{\GateOne{tgm}}}
\def\ColTGT#1#2{\SYMC{tgt}{#1}{#2}{\GateOne{tgt}}}
\IfFontExistsTF{Times New Roman}
  {\def\ColTNR#1#2{\SYMC{tnr}{#1}{#2}{\GateOne{tnr}}}}
  {\def\ColTNR#1#2{}}

\newcount\MTboxnum
\newsavebox\MTglyphbox
\MTboxnum=\allocationnumber
\def\MTprepare#1#2#3{%
  \global\setbox\MTglyphbox=\hbox{$\displaystyle\MTbody{#2}$}%
  \directlua{MTscan("\luaescapestring{#1}", "\luaescapestring{#3}", \the\MTboxnum)}}
\def\MToutput{%
  \refstepcounter{mtn}%
  {\color[HTML]{CC9933}\box\MTglyphbox}}
\ExplSyntaxOn
\cs_new_protected:Npn \MTbox #1#2
  {
    \str_if_eq:VnTF \MTselected {#1}
      { { \fboxsep=1.5pt \fbox {#2} } }
      {#2}
  }
\ExplSyntaxOff

\newcommand\FH[2]{{\footnotesize\shortstack{#1\\#2}}}
\def\INPUT{\input{unicode-math-table.tex}}
\def\TABLE{%
\par\noindent
\setlength\tabcolsep{4pt}%
\begin{longtable}[l]{@{}l *{\Nfonts}{c} lp{100pt}@{}}
  \toprule
  \textsc{usv} &
  \FH{XITS}{Math} &
  \FH{TeX Gyre}{Termes Math} &
  \FH{TeX Gyre}{Termes} &
  \FH{Times New}{Roman} &
  \FH{Times}{Math} &
  Macro & Description \\
  \midrule \endhead
  \INPUT\\
  \bottomrule
\end{longtable}
}
\makeatletter
\def\USV#1{\footnotesize\scshape\MakeLowercase{u+\@gobble#1}}
\makeatother
\def\CMD#1{\footnotesize\cmd#1}
\def\DESC#1{\vspace{-6pt}%
  \begin{spacing}{0.5}
    \RaggedRight\scriptsize #1%
  \end{spacing}
}

\newcommand\SHOW[1]{%
  \def\UnicodeMathSymbol##1##2##3##4{%
    \def\1{#1}\def\2{##3}%
    \ifx\1\2\PRINTLINE{##1}{##2}{##3}{##4}\fi
  }%
  \TABLE
}
\ExplSyntaxOn
\def\PRINTLINE#1#2#3#4{
    \MTprepare{#1}{#2}{#3}%
    \def\tempa{
      \USV{#1} &
      \MTbox{xits}{\ColXITS{#2}{#1}} &
      \MTbox{tgm}{\ColTGM{#2}{#1}} &
      \MTbox{tgt}{\ColTGT{#2}{#1}} &
      \MTbox{tnr}{\ColTNR{#2}{#1}} &
      \MTbox{mtn}{\MToutput} &
      \CMD{#2}
      \tl_if_in:NnT \PLAIN {#2}
        {
          \makebox[0pt][l]
            { \color[gray]{0.7} \textsuperscript{\sffamily (p)} }
        }
      \tl_if_in:NnT \LTXSYM {#2}
        {
          \makebox[0pt][l]
            { \color[gray]{0.7} \textsuperscript{\sffamily (l)} }
        }
      \tl_if_in:NnT \AMSSYMB {#2}
        {
          \makebox[0pt][l]
            { \color[gray]{0.7} \textsuperscript{\sffamily (a)} }
        }
      &
      \DESC{#4} \\
    }%
    \expandafter\tempa
}
\ExplSyntaxOff

\def\PLAIN{\alpha\beta\gamma\delta\epsilon\zeta\eta\theta\iota\kappa\lambda\mu\nu\xi\pi\rho\sigma\tau\upsilon\phi\chi\psi\omega\varepsilon\vartheta\varpi\varrho\varsigma\varphi\Gamma\Delta\Theta\Lambda\Xi\Pi\Sigma\Upsilon\Phi\Psi\Omega
%
\aleph\hbar\imath\jmath\ell\wp\Re\Im\partial\infty\prime\emptyset\nabla\surd\top\bot\angle\triangle\forall\exists\neg\flat\natural\sharp\clubsuit\diamondsuit\heartsuit\spadesuit
%
\coprod\bigvee\bigwedge\biguplus\bigcap\bigcup\int\prod\sum\bigotimes\bigoplus\bigodot\oint\bigsqcup\smallint
%
\triangleleft\triangleright\bigtriangleup\bigtriangledown\wedge\land\vee\lor\cap\cup\ddagger\dagger\sqcap\sqcup\uplus\amalg\diamond\bullet\wr\div\odot\oslash\otimes\ominus\oplus\mp\pm\circ\bigcirc\setminus\cdot\ast\times\star\propto\sqsubseteq\sqsupseteq\parallel\mid\dashv\vdash\nearrow\searrow\nwarrow\swarrow\Leftrightarrow\Leftarrow\Rightarrow\neq\ne\lnot\leq\le\geq\ge\succ\prec\approx\succeq\preceq\supset\subset\supseteq\subseteq\in\ni\owns\gg\ll\not\leftrightarrow\leftarrow\gets\rightarrow\to\mapstochar\mapsto\sim\simeq\perp\equiv\asymp\smile\frown\leftharpoonup\leftharpoondown\rightharpoonup\rightharpoondown
%
\joinrel\relbar\Relbar\lhook\hookrightarrow\rhook\hookleftarrow\bowtie\models\Longrightarrow\longrightarrow\longleftarrow\Longleftarrow\longmapsto\longleftrightarrow\Longleftrightarrow\iff
%
\ldotp\cdotp\colon\ldots\cdots\vdots\ddots
%
\acute\grave\ddot\tilde\bar\breve\check\hat\vec\dot\widetilde\widehat
%
\overrightarrow\overleftarrow\overbrace\underbrace\lmoustache\rmoustache\lgroup\rgroup\arrowvert\Arrowvert\bracevert\Vert\vert\uparrow\downarrow\updownarrow\Uparrow\Downarrow\Updownarrow\backslash\rangle\langle\rbrace\lbrace\rceil\lceil\rfloor\lfloor\sqrt}

\def\LTXSYM{
\cong
\notin
\rightleftharpoons
\doteq
\mathring
}

\def\AMSSYMB{\boxdot\boxplus\boxtimes\square\blacksquare\centerdot\lozenge\blacklozenge\circlearrowright\circlearrowleft\leftrightharpoons\boxminus\Vdash\Vvdash\vDash\twoheadrightarrow\twoheadleftarrow\leftleftarrows\rightrightarrows\upuparrows\downdownarrows\upharpoonright\restriction\downharpoonright\upharpoonleft\downharpoonleft\rightarrowtail\leftarrowtail\leftrightarrows\rightleftarrows\Lsh\Rsh\rightsquigarrow\leftrightsquigarrow\looparrowleft\looparrowright\circeq\succsim\gtrsim\gtrapprox\multimap\therefore\because\doteqdot\Doteq\triangleq\precsim\lesssim\lessapprox\eqslantless\eqslantgtr\curlyeqprec\curlyeqsucc\preccurlyeq\leqq\leqslant\lessgtr\backprime\risingdotseq\fallingdotseq\succcurlyeq\geqq\geqslant\gtrless\vartriangleright\vartriangleleft\trianglerighteq\trianglelefteq\bigstar\between\blacktriangledown\blacktriangleright\blacktriangleleft\vartriangle\blacktriangle\triangledown\eqcirc\lesseqgtr\gtreqless\lesseqqgtr\gtreqqless\Rrightarrow\Lleftarrow\veebar\barwedge\doublebarwedge\measuredangle\sphericalangle\varpropto\smallsmile\smallfrown\Subset\Supset\Cup\doublecup\Cap\doublecap\curlywedge\curlyvee\leftthreetimes\rightthreetimes\subseteqq\supseteqq\bumpeq\Bumpeq\lll\llless\ggg\gggtr\circledS\pitchfork\dotplus\backsim\backsimeq\complement\intercal\circledcirc\circledast\circleddash\lvertneqq\gvertneqq\nleq\ngeq\nless\ngtr\nprec\nsucc\lneqq\gneqq\nleqslant\ngeqslant\lneq\gneq\npreceq\nsucceq\precnsim\succnsim\lnsim\gnsim\nleqq\ngeqq\precneqq\succneqq\precnapprox\succnapprox\lnapprox\gnapprox\nsim\ncong\diagup\diagdown\varsubsetneq\varsupsetneq\nsubseteqq\nsupseteqq\subsetneqq\supsetneqq\varsubsetneqq\varsupsetneqq\subsetneq\supsetneq\nsubseteq\nsupseteq\nparallel\nmid\nshortmid\nshortparallel\nvdash\nVdash\nvDash\nVDash\ntrianglerighteq\ntrianglelefteq\ntriangleleft\ntriangleright\nleftarrow\nrightarrow\nLeftarrow\nRightarrow\nLeftrightarrow\nleftrightarrow\divideontimes\varnothing\nexists\Finv\Game\eth\eqsim\beth\gimel\daleth\lessdot\gtrdot\ltimes\rtimes\shortmid\shortparallel\smallsetminus\thicksim\thickapprox\approxeq\succapprox\precapprox\curvearrowleft\curvearrowright\digamma\varkappa\Bbbk\hslash\backepsilon}

\directlua{
  local function mtkey(n)
    local f = font.getfont(n.font)
    if not f then return nil, 0 end
    local name = string.lower(tostring(f.filename))
    local key = nil
    if string.find(name, "times-math-integrals", 1, true) then key = "mtn"
    elseif string.find(name, "times new roman", 1, true) then key = "tnr"
    elseif string.find(name, "texgyretermes-math", 1, true) then key = "tgm"
    elseif string.find(name, "texgyretermes", 1, true) then key = "tgt"
    elseif string.find(name, "xitsmath", 1, true) then key = "xits" end
    return key, f.size or 0
  end
  function MTscan(target, class, boxnum)
    local want = tonumber(string.gsub(target, '"', ''), 16)
    local exact, best, bestsize = nil, nil, -1
    local function walk(head)
      local n = head
      while n do
        if n.id == node.id("glyph") then
          local key, size = mtkey(n)
          if key then
            if want and n.char == want and not exact then exact = key end
            if size > bestsize then bestsize = size; best = key end
          end
        elseif n.id == node.id("hlist") or n.id == node.id("vlist") then
          walk(n.head)
        end
        n = n.next
      end
    end
    walk(tex.box[boxnum])
    local sel = exact
    if sel == nil then sel = best end
    token.set_macro("MTselected", sel or "", "global")
  end
}

\begin{document}
\MakeShortVerb\|
\title{Times Math\\[.3ex]\large Glyph reference}
\author{Junhao Wu}
\date{2026-10-17 \quad v0.1a}
\maketitle

This reference reads \texttt{unicode-math-table.tex} and groups its entries
by mathematical class and alphabet. Each row gives the Unicode scalar
value (USV), font samples, command name and description.

\begin{center}
\begin{tabular}{@{}lr@{}}
\toprule
Font or setup & Entries shown\\
\midrule
XITS Math, default features & \ref{count:xits}\\
TeX Gyre Termes Math & \ref{count:tgm}\\
TeX Gyre Termes, raw & \ref{count:tgt}\\
Times New Roman, raw & \ref{count:tnr}\\
\texttt{times-math} package output & \ref{count:mtn}\\
\bottomrule
\end{tabular}
\end{center}

The first four columns require the listed codepoint in the font's cmap;
an unavailable glyph leaves a blank cell. Counts refer to shown table
entries, including repeated codepoints. The Times Math column uses the
package's font ranges and alphabet mappings, including the bundled integral
font. If Times New Roman is unavailable, that column uses the mathematical
fonts.

A frame marks the font that supplies the glyph in the Times Math column.
The sixteen supplementary integrals have no source column, so their frames
appear in the Times Math column itself.

For comparison, the Times Math column uses \texttt{math-style=literal} and
\texttt{bold-style=literal}. Ordinary documents follow their own style
settings. The markers {\color[gray]{0.6}\textsuperscript{\sffamily (p)}},
{\color[gray]{0.6}\textsuperscript{\sffamily (l)}} and
{\color[gray]{0.6}\textsuperscript{\sffamily (a)}} identify command names
also used by Plain \TeX, \LaTeX\ and \textsf{amssymb}, respectively.

Accents use a common dotted-circle base (U+25CC) from XITS Math.
The TeX Gyre Termes and Times New Roman columns show text-font marks with a
single circle. The XITS Math, TeX Gyre Termes Math and Times Math columns
show mathematical attachment and stretching, using three circles for wide
examples.

This listing is derived from \textsf{unicode-math}'s
\texttt{unimath-symbols} document (Will Robertson et al.) as adapted by
\texttt{euler-math}, using per-font glyph checks and counts.

\tableofcontents

\clearpage
\section{Opening symbols, \cmd\mathopen}
\begingroup
\def\sqrt{\sqrtsign{}}
\def\cuberoot{\cuberootsign{}}
\def\fourthroot{\fourthrootsign{}}
\def\longdivision{\longdivisionsign{}}
\SHOW\mathopen
\endgroup

\section{Closing symbols, \cmd\mathclose}
\SHOW\mathclose

\section{Fence symbols, \cmd\mathfence}
\SHOW\mathfence

\section{Punctuation symbols, \cmd\mathpunct}
\SHOW\mathpunct

\clearpage
\section{`Over' symbols, \cmd\mathover}
\begingroup
\def\MTbody#1{#1{\MTwidebase}}
\def\MTRawBody#1#2#3{\Umathaccent 0 \csname symMT#1\endcsname #2\relax{{}\MTwidebase}}
\SHOW\mathover
\endgroup

\section{`Under' symbols, \cmd\mathunder}
\begingroup
\def\MTbody#1{#1{\MTwidebase}}
\def\MTRawBody#1#2#3{\Umathaccent bottom 0 \csname symMT#1\endcsname #2\relax{{}\MTwidebase}}
\SHOW\mathunder
\endgroup

\section{Accents, \cmd\mathaccent}

The dotted circle marks the base position. The Times Math column uses the
ordinary top accents selected by \texttt{times-math}.

\begingroup
\def\MTbody#1{#1{\MTcircle}}
\def\MTRawBody#1#2#3{\Umathaccent fixed 0 \csname symMT#1\endcsname #2\relax{\MTcircle}}
\SHOW\mathaccent
\endgroup

\section{Accents wide, \cmd\mathaccentwide}
\begingroup
\def\MTbody#1{#1{\MTwidebase}}
\def\MTRawBody#1#2#3{\Umathaccent 0 \csname symMT#1\endcsname #2\relax{\MTwidebase}}
\SHOW\mathaccentwide
\endgroup

\clearpage
\section{Bottom accents, \cmd\mathbotaccent}
\begingroup
\def\MTbody#1{#1{\MTcircle}}
\def\MTRawBody#1#2#3{\Umathaccent bottom fixed 0 \csname symMT#1\endcsname #2\relax{\MTcircle}}
\SHOW\mathbotaccent
\endgroup

\section{Bottom accents wide, \cmd\mathbotaccentwide}
\begingroup
\def\MTbody#1{#1{\MTwidebase}}
\def\MTRawBody#1#2#3{\Umathaccent bottom 0 \csname symMT#1\endcsname #2\relax{\MTwidebase}}
\SHOW\mathbotaccentwide
\endgroup

\section{Overlay accents, \cmd\mathaccentoverlay}
\begingroup
\def\MTbody#1{#1{\MTcircle}}
\let\MTRawBody\MTRawOverlay
\def\MTRawText#1#2{%
  \hbox{\fontfamily{\csname MTfamily#1\endcsname}\selectfont
    \ooalign{\hfil\MTcircle\hfil\crcr\hfil\char#2\relax\hfil\crcr}}}
\SHOW\mathaccentoverlay
\endgroup

\section{Big operators, \cmd\mathop}

The math-font and Times Math columns include limits, following
\textsf{unicode-math}'s \cmd\nolimits\ settings. The text-font columns
show the source glyph without mathematical resizing or limits.
In the Times Math column, integrals absent from Termes come from the bundled
Times Math Integrals font, constructed from Termes outlines with native
MATH size variants and script metrics.
\begingroup
\def\MTbody#1{#1_0^1}
\let\MTRawBody\MTRawOperator
\SHOW\mathop
\endgroup

\section{Binary relations, \cmd\mathbin}
\SHOW\mathbin

\section{Ordinary symbols, \cmd\mathord}
\SHOW\mathord

\clearpage
\section{Relation symbols, \cmd\mathrel}
\SHOW\mathrel

\clearpage
\section{Alphabetical symbols, \cmd\mathalpha}

\ExplSyntaxOn
\cs_set:Npn \UnicodeMathSymbol #1#2#3#4 {
  \str_if_eq:nnT {\mathalpha} {#3} {
    \exp_not:n { \UnicodeMathSymbol {#1}{#2}{#3}{#4} }
  }
}
\CatchFileEdef
  \ALPHA{unicode-math-table.tex}
  {\char_set_catcode_space:N \ }


\def\INPUT{\ALPHA}
\cs_new:Npn \SLOTS #1#2 {
  \cs_set:Npn \UnicodeMathSymbol ##1##2##3##4 {
    \bool_lazy_and:nnT
      { \int_compare_p:n {##1 >= #1} }
      { \int_compare_p:n {##1 <= #2} }
      {
        \PRINTLINE{##1}{##2}{##3}{##4}
      }
  }
  \TABLE
  \cs_set:Npn \UnicodeMathSymbol ##1##2##3##4 {
    \bool_lazy_or:nnT
      { \int_compare_p:n {##1 > #2} }
      { \int_compare_p:n {##1 < #1} }
      {
        \exp_not:n { \UnicodeMathSymbol {##1}{##2}{##3}{##4} }
      }
  }
  \edef\ALPHA{\ALPHA}
}

\ExplSyntaxOff

\subsection{Normal weight}

\subsubsection{Upright Greek, uppercase}
\SLOTS{"00391}{"003A9}

\subsubsection{Upright Greek, lowercase}
\SLOTS{"003B1}{"003F5}

\subsubsection{Italic, Latin, uppercase}
\SLOTS{"1D434}{"1D44D}

\subsubsection{Italic, Latin, lowercase}
\SLOTS{"1D44E}{"1D467}

\subsubsection{Italic Greek, uppercase}
\SLOTS{"1D6E2}{"1D6FA}

\subsubsection{Italic Greek, lowercase}
\SLOTS{"1D6FC}{"1D71B}

\subsubsection{Script, Latin, uppercase}
\SLOTS{"1D49C}{"1D4B5}

\subsubsection{Script, Latin, lowercase}
\SLOTS{"1D4B6}{"1D4CF}

\subsubsection{Fraktur, Latin, uppercase}
\SLOTS{"1D504}{"1D51C}

\subsubsection{Fraktur, Latin, lowercase}
\SLOTS{"1D51E}{"1D537}

\subsubsection{Blackboard, Latin, uppercase}
\SLOTS{"1D538}{"1D550}

\subsubsection{Blackboard, Latin, lowercase}
\SLOTS{"1D552}{"1D56B}

\subsubsection{Sans serif, Latin, uppercase}
\SLOTS{"1D5A0}{"1D5B9}

\subsubsection{Sans serif, Latin, lowercase}
\SLOTS{"1D5BA}{"1D5D3}

\subsubsection{Italic sans serif, Latin, uppercase}
\SLOTS{"1D608}{"1D621}

\subsubsection{Italic sans serif, Latin, lowercase}
\SLOTS{"1D622}{"1D63B}

\subsubsection{Typewriter, Latin, uppercase}
\SLOTS{"1D670}{"1D689}

\subsubsection{Typewriter, Latin, lowercase}
\SLOTS{"1D68A}{"1D6A3}

\subsection{Bold}

\subsubsection{Bold, Latin, uppercase}
\SLOTS{"1D400}{"1D419}

\subsubsection{Bold, Latin, lowercase}
\SLOTS{"1D41A}{"1D433}

\subsubsection{Bold Greek, uppercase}
\SLOTS{"1D6A8}{"1D6C0}

\subsubsection{Bold Greek, lowercase}
\SLOTS{"1D6C2}{"1D6E1}

\subsubsection{Bold italic, Latin, uppercase}
\SLOTS{"1D468}{"1D481}

\subsubsection{Bold italic, Latin, lowercase}
\SLOTS{"1D482}{"1D49B}

\subsubsection{Bold italic Greek, uppercase}
\SLOTS{"1D71C}{"1D734}

\subsubsection{Bold italic Greek, lowercase}
\SLOTS{"1D736}{"1D755}

\subsubsection{Bold script, Latin, uppercase}
\SLOTS{"1D4D0}{"1D4E9}

\subsubsection{Bold script, Latin, lowercase}
\SLOTS{"1D4EA}{"1D503}

\subsubsection{Bold fraktur, Latin, uppercase}
\SLOTS{"1D56C}{"1D585}

\subsubsection{Bold fraktur, Latin, lowercase}
\SLOTS{"1D586}{"1D59F}

\subsubsection{Bold sans serif, Latin, uppercase}
\SLOTS{"1D5D4}{"1D5ED}

\subsubsection{Bold sans serif, Latin, lowercase}
\SLOTS{"1D5EE}{"1D607}

\subsubsection{Bold italic sans serif, Latin, uppercase}
\SLOTS{"1D63C}{"1D655}

\subsubsection{Bold italic sans serif, Latin, lowercase}
\SLOTS{"1D656}{"1D66F}

\subsubsection{Bold sans serif Greek, uppercase}
\SLOTS{"1D756}{"1D76E}

\subsubsection{Bold sans serif Greek, lowercase}
\SLOTS{"1D770}{"1D78F}

\subsubsection{Bold italic sans serif Greek, uppercase}
\SLOTS{"1D790}{"1D7A8}

\subsubsection{Bold italic sans serif Greek, lowercase}
\SLOTS{"1D7AA}{"1D7C9}

\subsection{Miscellaneous}
\def\UnicodeMathSymbol#1#2#3#4{\PRINTLINE{#1}{#2}{#3}{#4}}
\TABLE

\FINISH
\mbox{}
\end{document}

\endinput
